\documentclass[10pt,a4paper]{amsart}
\usepackage[margin=19mm]{geometry}
\usepackage{amsmath,amssymb,mathtools}
\usepackage[scr=rsfs,cal=euler]{mathalfa}
\usepackage{xfrac}
\usepackage[unicode,hidelinks,bookmarks=false]{hyperref}
\newcommand{\ie}{i.e.\ }
\newcommand{\cf}{cf.\ }
\newcommand{\resp}{respectively\ }
\newcommand{\Subsection}{\subsection}
\providecommand{\nobreakdash}{\nobreak\hbox{-}}
\setlength{\emergencystretch}{2em}
\multlinegap0pt
\usepackage{mathtools}
\usepackage{amssymb}
\newtheorem{lemma}{Lemma}
\def\Arg{\operatorname{Arg}}
\title{MME for Mating Correspondences}
\author{Vanessa Matus De La Parra}
\date{Source: arXiv:2310.14330v2 (CC BY 4.0)}
\begin{document}
\maketitle

\noindent\emph{Source and licence.} MME for Mating Correspondences. Version arXiv:2310.14330v2 (CC BY 4.0), distributed by arXiv under the Creative Commons Attribution 4.0 International licence, http://creativecommons.org/licenses/by/4.0/, as recorded in the arXiv licence metadata for this exact version and shown on the abstract page https://arxiv.org/abs/2310.14330v2. Full licence notice: \texttt{licenses/arxiv-2310.14330-CC-BY-4.0.txt}.\par\smallskip

\noindent\textit{Editorial scope.} Counted unit: the article's lemma radius (source0.tex lines 382--384). The article's complete original proof of it is delivered with it (source0.tex lines 386--420, one \texttt{proof} environment of 35 lines). No mathematics is written, repaired or re-derived: the statement and the proof are the article's own lines, in the article's own notation, and no retained line is substituted or re-flowed.  No further source-local context is needed: every cross-reference of every retained line resolves inside the two delivered spans.  Nothing was omitted from any retained span, so the package carries no omission record.  No retained line contains a citation, so the wrapper carries no bibliography and no bibliography entry.  No retained line contains a figure environment, an image-inclusion command or drawing code, so no image is delivered and the package declares no asset dependency.  Editorial formatting only, in addition: an AMS \texttt{amsart} preamble that restates the article's own declarations and macros used by the retained lines, namely the environments \texttt{lemma} on separate counters, together with the standard packages those lines need.  The retained environments print in this document as Lemma 1 in order of appearance, whereas the article prints the same items with the numbering of its own full document.  The complete native source of the article as distributed in the arXiv e-print for this exact version is bundled unchanged as \texttt{sources/148882/source0.tex}.
\begin{lemma}\label{radius}
Let $R$ be a rational map. For all $\delta>0$ small enough, there exists $r_R(\delta)>0$ such that $\lim\limits_{\delta\to 0}r_R(\delta)=0$ and $\operatorname{diam}(U)\leq r_R(\delta)$, for each component $U$ of $R^{-1}(B(x,\delta))$, for all $x\in\widehat{\mathbb{C}}$.
\end{lemma}

\begin{proof}
Put $p\in\widehat{\mathbb{C}}$. By the Local Normal Form Theorem, there exists $(U(p),\Phi_{1,p},\Phi_{2,p},m(p))$ be so that 
\begin{itemize}
\item $U(p)$ is a topological disk containing $p$,
\item $\Phi_{1,p}(U(p))=\mathbb{D}$,
\item $\Phi_{2,p}(R(U(p))=\mathbb{D}$,
\item $\Phi_{2,p}\circ R\circ \Phi_{1,p}^{-1}(z)=z^{m(p)}$, for some $m(p)\in\mathbb{Z}$.
\end{itemize}
Since $\bigcup_{p\in\widehat{\mathbb{C}}}R(U(p))=\widehat{\mathbb{C}}$ and $\widehat{\mathbb{C}}$ is compact, there exist $p_1,\cdot,p_N\in\widehat{\mathbb{C}}$ satisfying that $\bigcup_{j=1}^N R(U(p_j))=\widehat{\mathbb{C}}$.
\\

Let $f(z)=z^m$. We know that for every $A\subset\mathbb{C}$, $f^{-1}(A)=\bigcup_{k=1}^{m}g_k(A)$, where $g_k(re^{i\theta})=\zeta^kr^{\frac{1}{m}}e^{i\frac{\theta}{m}}$ for $1\leq k\leq m$. 
Let $w\in\mathbb{D}$ and let $\epsilon>0$ be so that $B(z,2\epsilon)\subset\mathbb{D}$. 
\\

If $w=0$, then $g_k(B(w,\epsilon))=B(0,\delta^{\frac{1}{m}})$. Therefore, each component of $f^{-1}(B(w,\delta))$ has diameter at most $w\delta^{\frac{1}{m}}$.
\\

If $0\in B(w,\epsilon)$, then $B(w,\epsilon)\subset B(w,2\epsilon)\subset\mathbb{D}$. By above, each component of $f^{-1}(B(w,\epsilon))$ has diameter at most $2(w\epsilon)^{\frac{1}{m}}$.
\\

If $0\notin B(w,\epsilon)$ and $w'\in B(w,\epsilon)$, write $w'=|w'|e^{i\theta}$ with $\theta\in (\Arg(w)-\pi,\Arg(w)+\pi)$, where $\Arg(w)$ is the principal branch of logarithm. Then $|w'|\in(|w|-\epsilon,|w|+\epsilon)$ and $\theta\in(\Arg(w)-\epsilon,\Arg(w)+\epsilon)$. Thus,
\begin{eqnarray*}
|g_k(w')-g_k(w)|&=&\big| |w'|^{\frac{1}{m}}e^{i\frac{\theta}{m}}-|w|^{\frac{1}{m}}e^{i\frac{\theta}{m}}+|w|^{\frac{1}{m}}e^{i\frac{\theta}{m}}-|w|^{\frac{1}{m}}e^{i\frac{\Arg(z)}{m}} \big|\\
&=&\big| (|w'|^{\frac{1}{m}}-|w|^{\frac{1}{m}})e^{i\frac{\theta}{m}}+|w|^{\frac{1}{m}}(e^{i\frac{\theta}{m}}-e^{i\frac{\Arg(w)}{m}}) \big|\\
&<&\epsilon^{\frac{1}{m}}+\frac{\pi}{m}\epsilon.
\end{eqnarray*}
Therefore, each component of $f^{-1}(B(w,\epsilon))$ has diameter at most $\epsilon^{\frac{1}{m}}+\frac{\pi}{m}\epsilon$.

Observe that for each $1\leq j\leq N$, both $\Phi_{2,p_j}$ and $\Phi_{1,p_j}^{-1}$ are Lipschitz. Let $L$ be the maximum Lipschitz constant among all $\Phi_{2,p_j}$, $1\leq j\leq N$, and $K$ be the maximum Lipschitz constant among all $\Phi_{1,p_j}^{-1}$, $1\leq j\leq N$. Denote by $\delta_0$ the Lebesgue number of the covering $\lbrace R(U(p_j))\rbrace_{j=1}^N$ and let $\delta<\frac{\delta_0}{2}$. Then there exists $1\leq j\leq N$ so that $B(z,\delta)\subset R(U(p_j))$. Thus,
$$\Phi_{2,p_j}(B(z,\delta))\subset B(\Phi_{2,p}(z),L\delta)\subset\mathbb{D}.$$
Then, each component of $f^{-1}(\Phi_{2,p_j}(B(z,\delta)))$ has diameter at most $\max\lbrace 2(2L\delta)^{\frac{1}{m}},(L\delta)^{\frac{1}{m}}+\frac{\pi}{m}(L\delta)\rbrace$. Finally, since $R^{-1}(B(z,\delta))=\Phi_{1,p_j}^{-1}f^{-1}\Phi_{2,p_j}(B(z,\delta))$, then each component of $R^{-1}(B(z,\delta))$ has diameter at most
$$r_R(\delta)\coloneqq\max\left\lbrace 2K(2L\delta)^{\frac{1}{m}},K((L\delta)^{\frac{1}{m}}+\dfrac{\pi}{m}(L\delta))\right\rbrace.$$
From the definition of $r_R(\delta)$, it is clear that $\lim\limits_{\delta\to 0}r_R(\delta)=0$. 
\end{proof}

\end{document}
